\section{Comparison with the known counterexamples}
\label{sec:counterexamples}

The uncorrected strong inequality fails at small volume, even on
Calabi--Yau threefolds.  We record the parameters of two established
counterexamples and then give a uniform estimate for sheaves on their
exceptional surfaces.  Throughout this section, the volume parameter is
\(a\) in the ordinary tilt slope \(\nu_{a,b}\); the physical volume is
\(t=\sqrt3a\).

\subsection{A line bundle and a contracted divisor}

Let \(X=\operatorname{Bl}_p\mathbb P^3\), let \(L\) be the pulled-back
hyperplane class, and let \(E\) be the exceptional divisor.  Fix
\(H=2L-E\) and \(B_0=0\).  The intersection identities
\(L^3=E^3=1\), \(LE=0\), and \(H^3=7\) give
\[
 \bigl(H^3\ch_0,H^2\ch_1,H\ch_2,\ch_3\bigr)(\OO_X(L))
 =\left(7,4,1,\frac16\right).
\]
For the conventional quadratic expression
\[
 \mathcal Q_{a,b}(F)
 =a^2\bigl((H^2\ch_1(F))^2-2H^3\ch_0(F)H\ch_2(F)\bigr)
  +4(H\ch_2^{bH}(F))^2
  -6H^2\ch_1^{bH}(F)\ch_3^{bH}(F),
\]
direct expansion yields
\begin{equation}
 \mathcal Q_{a,b}(\OO_X(L))
 =2\left(a^2+\left(b-\frac14\right)^2-\frac1{16}\right).
 \label{eq:schmidt-disk}
\end{equation}
Schmidt proves that actual tilt-stable points occur in the negative
region of \eqref{eq:schmidt-disk}
\cite[Theorem~3.1]{Schmidt2017}.  All these violations have \(a<1/4\).
On the positive-denominator zero-slope branch one has, more explicitly,
\[
 a^2=b^2-\frac87b+\frac27,
 \qquad b<\frac{4-\sqrt2}{7},
 \qquad
 \ch_3^{bH}(\OO_X(L))
 -\frac{a^2}{6}H^2\ch_1^{bH}(\OO_X(L))
 =\frac{4b-1}{42}.
\]
These formulas use the same ordinary characters and normalization as
the strong inequality.

For an integral effective divisor \(D\subset X\), still with \(B_0=0\),
the identity \(\ch(\OO_D)=1-e^{-D}\) shows that its unique zero-slope
line and its defect there are
\begin{align}
 b_0&=-\frac{HD^2}{2H^2D},\notag\\
 \ch_3^{b_0H}(\OO_D)
 -\frac{a^2}{6}H^2D
 &=\frac{D^3}{6}
   -\frac{(HD^2)^2}{8H^2D}
   -\frac{a^2}{6}H^2D.
 \label{eq:divisor-counterexample}
\end{align}
The wall-radius argument in
\cite[Lemma~3.1 and its proof]{MartinezSchmidt2019} gives semistability
on this line for \(a\ge H^2D/(2H^3)\), and stability for the strict
inequality in our subobject-versus-quotient convention.  Indeed, every
positive-rank wall has radius at most this bound, while a proper
rank-zero subobject of \(\OO_D\) is a subsheaf with quotient supported
in dimension at most one, hence of infinite tilt slope.

Suppose now that \(D\) contracts to a point, \(-D\) is relatively
ample, and \(H=ML-D\), where \(L\) is the pullback of an ample
divisor and \(M\) is sufficiently large.  Put \(s=D^3>0\).
Since \(LD=0\),
one has \(H^2D=s\) and \(HD^2=-s\).  Formula
\eqref{eq:divisor-counterexample} becomes
\begin{equation}
 b_0=\frac12,
 \qquad
 \ch_3^{b_0H}(\OO_D)-\frac{a^2}{6}H^2D
 =s\left(\frac1{24}-\frac{a^2}{6}\right).
 \label{eq:contracted-defect}
\end{equation}
Thus, when \(H^3>s\), the interval
\(s/(2H^3)<a<1/2\) consists of stable counterexamples.
The polarization, including \(M\), is fixed in this assertion.
Tensoring by \(\OO_X(kH)\) translates \(b\) by \(k\) and preserves
the twisted characters and \(a\); it gives counterexamples with
unbounded \(b\), but with the same bounded volume interval.

\subsection{A uniform estimate on the reduced exceptional surface}

The same upper bound on the violating volume applies to sheaves of
arbitrary rank on a smooth reduced exceptional surface, whenever their
pushforwards are tilt-semistable.  The following statement also allows
twists transverse to the polarization.

\begin{lemma}
\label{lem:exceptional-surface-bound}
Let \(X\) be a smooth projective complex threefold, let \(H\) be an
ample integral divisor, and let \(i:D\hookrightarrow X\) be a smooth
integral divisor such that \(H|_D=-D|_D\).  Fix
\(B_0\in N^1(X)_{\QQ}\), and set \(B=B_0+bH\).
If \(F\) is a torsion-free sheaf of positive rank on \(D\) and
\(i_*F\) is \(\nu_{a,b}\)-semistable of slope zero, then
\begin{equation}
 \frac{\ch_3^B(i_*F)}{H^2\ch_1^B(i_*F)}\le\frac1{24}.
 \label{eq:exceptional-uniform}
\end{equation}
Consequently these objects satisfy the corrected strong inequality
with \(\Gamma=H^2/24\) for every \(a>0\), and the uncorrected strong
inequality for every \(a\ge1/2\).  Both constants are independent of
\(F\), its rank, and \(b\).
\end{lemma}

\begin{proof}
Write \(\ell=H|_D\), \(s=\ell^2=D^3>0\),
\(\rho=\rk(F)\), and \(C=B_0|_D\).  Put
\[
 u=\ch_1^C(F),\qquad v=\ch_2^C(F),\qquad
 \mu=\frac{\ell u}{\rho s},\qquad
 \delta=\mu^2-\frac{2v}{\rho s}.
\]
Every exact sequence of sheaves on \(D\) pushes forward to an exact
sequence of torsion sheaves in the first-tilt heart.  Their tilt
slopes differ from their \(C\)-twisted normalized surface slopes by
the common constant \(1/2-b\).  Twisting shifts all surface slopes
equally, so tilt semistability of \(i_*F\) implies slope semistability
of \(F\) for \(\ell\).
The classical surface Bogomolov--Gieseker inequality
\cite[Theorem~3.1.4]{BMT2014} and the Hodge index theorem give
\[
 u^2-2\rho v\ge0,
 \qquad (\ell u)^2\ge s u^2,
 \qquad \delta\ge0.
\]
The first expression is unchanged by twisting the Chern character
by \(C\), so these conclusions hold for the specified \(B_0\).

The normal line bundle has first Chern class \(-\ell\), and hence
\(\td(N_{D/X})^{-1}=1+\ell/2+\ell^2/6\).
Grothendieck--Riemann--Roch now gives
\[
 H^2\ch_1^B(i_*F)=\rho s,
 \qquad
 H\ch_2^B(i_*F)=\rho s(\mu+1/2-b).
\]
Since the threefold rank of \(i_*F\) is zero, its zero-slope condition
is \(b=\mu+1/2\).  The degree-three term of the same formula yields
\[
 \frac{\ch_3^B(i_*F)}{\rho s}
 =\frac{v}{\rho s}+\frac\mu2+\frac16
   -b(\mu+1/2)+\frac{b^2}{2}
 =\frac1{24}-\frac\delta2\le\frac1{24}.
\]
This proves \eqref{eq:exceptional-uniform}; the two strong bounds
follow by subtracting \(a^2/6\).
\end{proof}

For the Weierstrass examples of
\cite[Section~3.2 and Appendix~A]{MartinezSchmidt2019}, let
\(p:X\to S\) be a Weierstrass elliptic fibration from a smooth
Calabi--Yau threefold to a del Pezzo surface, and let \(\Sigma\)
be its section.
Since \(K_X\simeq\OO_X\), adjunction gives
\(\Sigma|_\Sigma=K_S\).  For the ample polarization
\(H=q\Sigma-(1+q)p^*K_S\), with a positive integer \(q\), one has
\(H|_\Sigma=-K_S\).  Thus
Lemma~\ref{lem:exceptional-surface-bound} applies in arbitrary rank
and after every line-bundle twist on \(\Sigma\) whose pushforward is
tilt-semistable.  In particular,
the known Calabi--Yau counterexample has the defect
\eqref{eq:contracted-defect}, with \(s=K_S^2\), and its violating
volume range is bounded by \(a<1/2\).

\begin{remark}
\label{rem:nonreduced-multiple}
Replacing \(D\) by \(nD\) in the character calculation does not
produce a stable counterexample at arbitrarily large volume.  Under
the hypotheses of Lemma~\ref{lem:exceptional-surface-bound}, the exact
sequence
\[
 0\longrightarrow\OO_D(-(n-1)D)\longrightarrow\OO_{nD}
 \longrightarrow\OO_{(n-1)D}\longrightarrow0
\]
lies in every first-tilt heart.  For \(B_0=0\) and \(b=n/2\), its
three slopes are \((n-1)/2\), \(0\), and \(-1/2\), respectively.
Thus \(\OO_{nD}\) is unstable for every \(a>0\) when \(n>1\).
The uniform estimate above concerns sheaves on the reduced divisor;
its proof uses surface semistability, rather than the character of
an arbitrary multiple.
\end{remark}
